\section*{NeurIPS Paper Checklist}

\begin{enumerate}[leftmargin=*]

\item \textbf{Claims}\\
Question: Do the main claims made in the abstract and introduction accurately reflect the paper's
contributions and scope?\\
Answer: [Yes]\\
Justification: The abstract and Sections 1--3 state the paper's main contribution as executable
manipulation checks for role-prompted multi-agent systems, distinguish observable behavioral
signatures from latent computation, and report the post-collection guard audit. Section 6 explicitly
limits generalization beyond the two model families and benchmarks studied.

\item \textbf{Limitations}\\
Question: Does the paper discuss the limitations of the work performed by the authors?\\
Answer: [Yes]\\
Justification: Section 6 provides a dedicated discussion of limitations in scope, measurement,
provenance, statistical power, and departures from the preregistered design. Appendix D
additionally reports limitations discovered by post-collection audit, including collapsed or
defeatable manipulation-check guards.

\item \textbf{Theory assumptions and proofs}\\
Question: For each theoretical result, does the paper provide the full set of assumptions and a
complete (and correct) proof?\\
Answer: [N/A]\\
Justification: The paper introduces an operational definition of an executable manipulation check,
but does not claim new theorems, lemmas, or formal theoretical results requiring proofs.

\item \textbf{Experimental result reproducibility}\\
Question: Does the paper fully disclose all the information needed to reproduce the main
experimental results of the paper to the extent that it affects the main claims and/or conclusions of
the paper (regardless of whether the code and data are provided or not)?\\
Answer: [Yes]\\
Justification: Sections 3--4 specify the experimental conditions, decision rules, models, decoding
settings, benchmark revisions, and attribution procedure. Appendices E, I, J, and O provide the split
construction, exact manipulation-check definitions, robustness analyses, and rendered prompts needed
to reproduce the main comparisons.

\item \textbf{Open access to data and code}\\
Question: Does the paper provide open access to the data and code, with sufficient instructions to
faithfully reproduce the main experimental results, as described in supplemental material?\\
Answer: [Yes]\\
Justification: The full artifact is available at
\url{https://github.com/kevzho/role-prompted-systems} under an MIT license with
complete environment and reproduction commands. The paper identifies the public benchmark revisions
and documents the experimental procedure in detail.

\item \textbf{Experimental setting/details}\\
Question: Does the paper specify all the training and test details (e.g., data splits,
hyperparameters, how they were chosen, type of optimizer) necessary to understand the results?\\
Answer: [Yes]\\
Justification: Section 4 specifies the benchmark arms, model deployments, temperature, token limit,
seed, timeout, preregistered decision thresholds, and attribution estimator. Appendix E specifies the
hash-based train/calibration/validation/test partition, and Appendix O reproduces the prompts used in
the manipulation-check experiment.

\item \textbf{Experiment statistical significance}\\
Question: Does the paper report error bars suitably and correctly defined or other appropriate
information about the statistical significance of the experiments?\\
Answer: [Yes]\\
Justification: Section 3.4 and Section 4.1 specify paired item-level bootstrap confidence intervals
using 2,000 resamples, Benjamini--Hochberg correction for the registered role-check family, and
preregistered effect-size thresholds. The results report confidence intervals, corrected significance
values, and effect sizes, with additional sensitivity and resampling analyses in Appendix J.

\item \textbf{Experiments compute resources}\\
Question: For each experiment, does the paper provide sufficient information on the computer
resources (type of compute workers, memory, time of execution) needed to reproduce the experiments?\\
Answer: [Yes]\\
Justification: Section 4.3 identifies the hosted model provider, deployment route, decoding
parameters, and model versions, provider-side hardware, memory, execution time, call counts, and
token usage.

\item \textbf{Code of ethics}\\
Question: Does the research conducted in the paper conform, in every respect, with the NeurIPS Code
of Ethics \url{https://neurips.cc/public/EthicsGuidelines}?\\
Answer: [Yes]\\
Justification: The study evaluates existing hosted language models on established benchmarks, does
not involve human-subject interventions, and makes no deployment claims. Section 6 explicitly
cautions against using the evaluated council for medical, legal, financial, or other consequential
decisions.

\item \textbf{Broader impacts}\\
Question: Does the paper discuss both potential positive societal impacts and negative societal
impacts of the work performed?\\
Answer: [No]\\
Justification: The paper discusses risks of over-interpreting benchmark performance and explicitly
cautions against consequential deployment in Section 6, but the current manuscript does not contain a
dedicated discussion of both positive and negative societal impacts of executable manipulation
checks.

\item \textbf{Safeguards}\\
Question: Does the paper describe safeguards that have been put in place for responsible release of
data or models that have a high risk for misuse (e.g., pre-trained language models, image generators,
or scraped datasets)?\\
Answer: [N/A]\\
Justification: The paper does not release a new pretrained model, high-risk dataset, or other asset
requiring controlled access. It evaluates existing hosted models and established benchmarks.

\item \textbf{Licenses for existing assets}\\
Question: Are the creators or original owners of assets (e.g., code, data, models), used in the
paper, properly credited and are the license and terms of use explicitly mentioned and properly
respected?\\
Answer: [Yes]\\
Justification: The paper cites the original HaluEval and TruthfulQA publications and identifies
pinned benchmark revisions; the released artifact redistributes both corpora under their original
licenses, with attribution in LICENSE.

\item \textbf{New assets}\\
Question: Are new assets introduced in the paper well documented and is the documentation provided
alongside the assets?\\
Answer: [N/A]\\
Justification: The current submission does not release a new dataset, pretrained model, or separately
distributed software asset. The new experimental prompts and manipulation-check specifications are
documented directly in the paper and appendices.

\item \textbf{Crowdsourcing and research with human subjects}\\
Question: For crowdsourcing experiments and research with human subjects, does the paper include the
full text of instructions given to participants and screenshots, if applicable, as well as details
about compensation (if any)?\\
Answer: [N/A]\\
Justification: The study does not involve crowdsourcing or human subjects.

\item \textbf{Institutional review board (IRB) approvals or equivalent for research with human
subjects}\\
Question: Does the paper describe potential risks incurred by study participants, whether such risks
were disclosed to the subjects, and whether Institutional Review Board (IRB) approvals (or an
equivalent approval/review based on the requirements of your country or institution) were obtained?\\
Answer: [N/A]\\
Justification: The study does not involve crowdsourcing or human subjects, so IRB or equivalent
human-subjects review is not applicable.

\item \textbf{Declaration of LLM usage}\\
Question: Does the paper describe the usage of LLMs if it is an important, original, or non-standard
component of the core methods in this research? Note that if the LLM is used only for writing,
editing, or formatting purposes and does not impact the core methodology, scientific rigor, or
originality of the research, declaration is not required.\\
Answer: [Yes]\\
Justification: LLMs are a core component of the experiment, while also being a writing and coding
aid. Section 4.3 identifies the model families, hosted serving environment, decoding parameters, and
model provenance, while Appendix O reproduces the role prompts and shared output contract used to
generate the experimental records.

\end{enumerate}
